\par\Needspace{6\baselineskip}\medskip
\hypertarget{native-142-article103}{}\pdfbookmark[2]{103.}{article-103}
\begin{center}103.\end{center}
Quoties $z$ valorem maiorem obtinet, tabulae III limites egredientem, aequationes 13, $13^*$ tuto semper ac commode in forma sua non mutata tentando solventur, et quidem ob rationes iis similes quas in art. 94 pro ellipsi exposuimus. In tali casu elementa orbitae obiter saltem cognita esse supponere \SourceLeaf{143}licet: tum vero valor approximatus ipsius $n$ statim habetur per formulam $\operatorname{tang}2n=\dfrac{\sin f\mathpunct{.}\BareRoot{rr'}}{\alpha\BareRoot{(ee-1)}}$, quae sponte demanat ex aequatione 6 art. 99. Ex $n$ autem habebitur $z$ per formulam $z=\dfrac{1-\cos2n}{2\cos2n}=\dfrac{\sin n^2}{\cos2n}$, et ex valore approximato ipsius $z$ paucis tentaminibus derivabitur ille, qui aequationi 13 vel $13^*$ ex asse satisfacit. Possunt quoque illae aequationes in hac forma exhiberi
\[
\begin{aligned}
m={}&\left(l-\frac{\sin n^2}{\cos2n}\right)^{\frac12}\\[.6em]
&+2\left(l-\frac{\sin n^2}{\cos2n}\right)^{\frac32}\frac{\dfrac{\operatorname{tang}2n}{\cos2n}-\log\operatorname{hyp}\operatorname{tang}(45^0+n)}{\operatorname{tang}2n^3}
\end{aligned}
\]
\[
\begin{aligned}
M={}&-\left(L+\frac{\sin n^2}{\cos2n}\right)^{\frac12}\\[.6em]
&+2\left(L+\frac{\sin n^2}{\cos2n}\right)^{\frac32}\frac{\dfrac{\operatorname{tang}2n}{\cos2n}-\log\operatorname{hyp}\operatorname{tang}(45^0+n)}{\operatorname{tang}2n^3}
\end{aligned}
\]
\noindent atque sic, neglecta $z$, statim valor verus ipsius $n$ erui.
